\begin{frame}{nDCG — interprétation du gain linéaire vs exponentiel}

%\textbf{Formule générale :}
\[
\text{DCG@k} = \sum_{i=1}^k \frac{g_i}{\log_2(i+1)} \qquad 
\text{où } g_i = \text{gain}(rel_i)
\]

\begin{block}{Deux choix possibles pour la fonction de gain}
\begin{itemize}
    \item \textbf{Gain linéaire} : $g(rel) = rel$  
    $\Rightarrow$ chaque niveau de pertinence ajoute une valeur proportionnelle.
    \item \textbf{Gain exponentiel} : $g(rel) = 2^{rel}-1$  
    $\Rightarrow$ les documents très pertinents comptent beaucoup plus.
\end{itemize}
\end{block}

\begin{center}
\begin{tabular}{ccccc}
\toprule
Niveau de pertinence & 0 & 1 & 2 & 3 \\
\midrule
Gain linéaire & 0 & 1 & 2 & 3 \\
Gain exponentiel & 0 & 1 & 3 & 7 \\
\bottomrule
\end{tabular}
\end{center}
\end{frame}

\begin{frame}{nDCG — interprétation du gain linéaire vs exponentiel}
\begin{center}
\begin{tabular}{ccccc}
\toprule
Niveau de pertinence & 0 & 1 & 2 & 3 \\
\midrule
Gain linéaire & 0 & 1 & 2 & 3 \\
Gain exponentiel & 0 & 1 & 3 & 7 \\
\bottomrule
\end{tabular}
\end{center}
\vspace{0.5em}
\textbf{Interprétation :}
\begin{itemize}
  \item Le \textbf{gain linéaire} favorise la régularité : plusieurs résultats pertinents suffisent.
  \item Le \textbf{gain exponentiel} surpondère les très pertinents : utile pour les classements courts (web search).
\end{itemize}

\end{frame}
